\section{Convergence diagnostics and auditable terminal states}
The October 2026 supplementary campaign uses the shared task interface and retains execution status, physical acceptance, source identity and input identity separately. The manufactured solution passed on 208, 756 and 1856 nodes; the largest field and derivative errors were $1.70\times10^{-14}$ and $2.26\times10^{-13}$.

Analytical boundary traces and traces from the discrete two-dimensional solution were recomputed on paired extruded meshes with 1683, 3672 and 9108 nodes. On the finest mesh, the analytical-trace current error was 0.877248\%, the $\phi$ error was $1.70939\times10^{-4}$, and the gauge residual was $2.77686\times10^{-2}$. The discrete-trace values were 0.877132\%, $1.47596\times10^{-4}$ and $2.77850\times10^{-2}$. Both controls retain the computed $\phi$ and its derivatives. Replacing the boundary trace alone did not satisfy the $\phi$ and gauge limits of $10^{-8}$ and $10^{-3}$. The next minimum experiment is a block-residual test of the extruded discrete operator and its gauge coupling.

The seven-strand Maxwell case used the official mesh-link interface and one evidence-based settings revision. Source and basis designs both contained 1694 actual mesh elements, but exact coordinate identity still failed. Canonicalization permits node and element renumbering only. The retained outcome is a fixed-mesh diagnostic; spatial convergence, boundary sensitivity and a matching cross-solver comparison remain gated. Historical single-mesh agreement retains its original scope.

The machine-readable evidence index is in the repository's provenance/convergence directory. These finite outcomes guide subsequent experiments. Final Q2 design still requires at least $5\,\mathrm{mm^2}$ nominal copper area, a traditional multilevel Litz baseline, and fair actual-length and copper-volume accounting.
